\newcommand{\figuraPasajeroAutomovil}{%
\begin{figure}[H]
    \centering
    \begin{tikzpicture}[>=Latex,font=\small,line cap=round,line join=round]
        \fill[green!8,rounded corners=4pt] (-0.5,-0.5) rectangle (10.5,5.2);

        \draw[line width=15mm,gray!38]
            (0.2,0.15) .. controls (0.0,2.7) and (2.4,4.7) .. (6.2,4.75);
        \draw[white,line width=1pt,dashed]
            (0.2,0.15) .. controls (0.0,2.7) and (2.4,4.7) .. (6.2,4.75);

        \coordinate (P) at (2.25,3.72);
        \fill[red!75!black] (P) circle (0.14);
        \node[anchor=north west] at ($(P)+(0.10,-0.12)$) {automóvil};

        \draw[->,very thick,dashed,black!65]
            (P)--++(2.65,0.74);
        \draw[->,ultra thick,green!45!black]
            (P)--++(-1.25,0.95);
        \draw[->,ultra thick,red!75!black]
            (P)--++(1.15,-1.05);

        \fill[blue!70!black] (0.45,4.85) circle (0.11);
        \node[anchor=south] at (0.45,5.0) {centro de curvatura};
        \draw[densely dashed,blue!60!black] (0.45,4.85)--(P);

        \draw[->,ultra thick,green!45!black] (6.7,4.0)--(7.6,4.0);
        \node[anchor=west] at (7.75,4.0) {fuerza real hacia el centro};
        \draw[->,very thick,dashed,black!65] (6.7,3.0)--(7.6,3.0);
        \node[anchor=west] at (7.75,3.0) {dirección tangente};
        \draw[->,ultra thick,red!75!black] (6.7,2.0)--(7.6,2.0);
        \node[anchor=west] at (7.75,2.0) {fuerza ficticia hacia afuera};
        \node[align=left,anchor=west,text=black!70] at (6.7,0.85)
            {La flecha roja solo se usa\\en el marco del automóvil.};
    \end{tikzpicture}
    \caption{Representación simplificada del automóvil como un punto sobre
    una trayectoria curva. En el marco inercial, la fuerza real modifica la
    dirección del movimiento; en el marco del automóvil aparece una fuerza
    ficticia en sentido opuesto.}
    \label{fig:pasajero-automovil}
\end{figure}%
}